\chapter{Referencia rápida de comandos}
\label{app:comandos}

Este apéndice reúne los comandos principales del procedimiento en Linux.
Las explicaciones y condiciones de uso se encuentran en los capítulos correspondientes.

\section{Comprobar el sistema}

\begin{lstlisting}[language=bash]
uname -r
uname -m
cat /etc/os-release
lspci -nn
\end{lstlisting}

\section{Raspberry Pi 5}

Para un adaptador PCIe no HAT+, edita:

\begin{lstlisting}[language=bash]
sudoedit /boot/firmware/config.txt
\end{lstlisting}

y añade, si todavía no está presente:

\begin{lstlisting}
dtparam=pciex1
\end{lstlisting}

Después reinicia y comprueba la conexión:

\begin{lstlisting}[language=bash]
sudo reboot
lspci -nn
\end{lstlisting}

\section{Controlador PCIe}

En Ubuntu:

\begin{lstlisting}[language=bash]
sudo apt update
sudo apt install dkms build-essential "linux-headers-$(uname -r)"
\end{lstlisting}

Descarga e instala el controlador:

\begin{lstlisting}[language=bash]
git clone https://github.com/Brainchip-Inc/akida_dw_edma.git
cd akida_dw_edma
git rev-parse HEAD
./install.sh
\end{lstlisting}

Comprueba la instalación:

\begin{lstlisting}[language=bash]
dkms status
lsmod | grep -i akida
ls -l /dev/akida*
lspci -nnk
\end{lstlisting}

\section{Entorno de Python}

Desde la carpeta principal del proyecto:

\begin{lstlisting}[language=bash]
python3.11 -m venv .venv
source .venv/bin/activate
python -m pip install --upgrade pip
\end{lstlisting}

Instalación actual de MetaTF:

\begin{lstlisting}[language=bash]
python -m pip install "metatf==2.19.3"
\end{lstlisting}

o entorno de referencia del repositorio:

\begin{lstlisting}[language=bash]
python -m pip install -r requirements/reproducible.txt
\end{lstlisting}

Instala después las herramientas del proyecto:

\begin{lstlisting}[language=bash]
python -m pip install -e . --no-deps
python -m pip check
\end{lstlisting}

\section{Pruebas básicas}

Prueba en software:

\begin{lstlisting}[language=bash]
python examples/00_simulator_smoke.py
akd1000-bringup verify \
  --id smoke \
  --model models/akida_v1_smoke_test.fbz
\end{lstlisting}

Diagnóstico:

\begin{lstlisting}[language=bash]
akd1000-bringup doctor
\end{lstlisting}

Prueba en la tarjeta:

\begin{lstlisting}[language=bash]
akida devices
python examples/01_hardware_smoke.py
\end{lstlisting}

\section{Ejecutar un modelo propio}

\begin{lstlisting}[language=bash]
akd1000-bringup run \
  --model models/model.fbz \
  --input data/input.npy \
  --method forward \
  --clock Performance \
  --map AllNps \
  --output outputs/raw-output.npy
\end{lstlisting}

\section{Benchmark}

\begin{lstlisting}[language=bash]
akd1000-bringup benchmark \
  --model models/model.fbz \
  --input data/inputs.npy \
  --method forward \
  --clock Performance \
  --map AllNps \
  --warmup 8 \
  --repetitions 1 \
  --output outputs/smoke
\end{lstlisting}

\section{Diagnóstico rápido}

\begin{enumerate}
  \item Si la tarjeta no aparece en \code{lspci}, revisa el montaje y la configuración PCIe.
  \item Si aparece en PCIe pero falta \code{/dev/akida*}, revisa el controlador.
  \item Si existe el nodo pero \code{akida.devices()} está vacío, revisa el entorno y los permisos.
  \item Si el modelo de prueba funciona y el modelo propio no, revisa el modelo y sus datos de entrada.
\end{enumerate}